\documentclass[../../main.tex]{subfiles}

\begin{document}

\section{Baire 范畴定理}\label{section:ch10-baire}

令 $E$ 为度量空间 $X$ 的子集. 点 $x\in E$ 称为 $E$ 的内点, 若存在包含于 $E$ 的、中心在 $x$ 的开球. $E$ 的内点全体称为 $E$ 的内部且记为 $\operatorname{int}E$. 点 $x\in X\sim E$ 称为 $E$ 的外点, 若存在包含于 $X\sim E$ 的、中心在 $x$ 的开球. $E$ 的外点全体称为 $E$ 的外部且记为 $\operatorname{ext}E$. 若点 $x\in X$ 具有每个中心在 $x$ 的球包含 $E$ 与 $X\sim E$ 中的点的性质, 则称为 $E$ 的边界点; $E$ 的边界点全体称为 $E$ 的边界且记为 $\operatorname{bd}E$. 我们把证明对 $X$ 的任何子集 $E$
\begin{align}
X=\operatorname{int}E\cup\operatorname{ext}E\cup\operatorname{bd}E\text{ 且这个并是不交的}\label{align:ch10-baire-5}
\end{align}
留作练习.

回忆度量空间 $X$ 的子集 $A$ 称为(在 $X$ 中)稠密, 若 $X$ 的每个非空开子集包含 $A$ 中的点. 我们称度量空间的子集是(在 $X$ 中)中空的, 若它有空的内部. 观察到对于度量空间 $X$ 的子集 $E$,
\begin{align}
E\text{ 在 }X\text{ 中是中空的当且仅当它的补 }X\sim E\text{ 在 }X\text{ 中稠密}.\label{align:ch10-baire-6}
\end{align}
对于度量空间 $X$, 点 $x\in X$ 和 $0<r_{1}<r_{2}$, 由度量的连续性, 我们有包含关系 $\overline{B(x,r_{1})}\subseteq B(x,r_{2})$. 因此 $\overline{B(x,r_{1})}$ 是使得
\begin{align*}
B(x,r_{1})\subseteq\overline{B(x,r_{1})}\subseteq B(x,r_{2})
\end{align*}
的闭集. 因此, 若 $\mathcal{O}$ 是度量空间 $X$ 的开子集, 对每个点 $x\in\mathcal{O}$, 存在中心在 $x$ 的开球其闭包包含于 $\mathcal{O}$.

\begin{theorem}[Baire 范畴定理]\label{theorem:ch10-baire}
令 $X$ 为完备度量空间.
\begin{enumerate}[(i)]
\item 令 $\{\mathcal{O}_{n}\}_{n=1}^{\infty}$ 为 $X$ 的开稠密子集的可数族. 则交 $\bigcap_{n=1}^{\infty}\mathcal{O}_{n}$ 也是稠密的.
\item 令 $\{F_{n}\}_{n=1}^{\infty}$ 为 $X$ 的闭中空子集的可数族. 则并 $\bigcup_{n=1}^{\infty}F_{n}$ 也是中空的.
\end{enumerate}
\end{theorem}

\begin{proof}
一个集合是稠密的当且仅当它的补是中空的. 一个集合是开的当且仅当它的补是闭的. 我们因此从 De Morgan 等式推出(i)和(ii)是等价的. 我们证明(i). 令 $x_{0}$ 属于 $X$ 而 $r_{0}>0$. 我们必须证明 $B(x_{0},r_{0})$ 包含 $\bigcap_{n=1}^{\infty}\mathcal{O}_{n}$ 中的点. 由于 $\mathcal{O}_{1}$ 在 $X$ 中稠密, 集合 $B(x_{0},r_{0})\cap\mathcal{O}_{1}$ 是非空的.

令 $x_{1}$ 属于开集 $B(x_{0},r_{0})\cap\mathcal{O}_{1}$. 选取 $r_{1}$, $0<r_{1}<1$, 使得若我们定义 $B_{1}=B(x_{1},r_{1})$, 则
\begin{align}
\overline{B}_{1}\subseteq B(x_{0},r_{0})\cap\mathcal{O}_{1}.\label{align:ch10-baire-7}
\end{align}
假定 $n$ 是自然数且开球的递减族 $\{B_{k}\}_{k=1}^{n}$ 已被选取. 它具有性质: 对 $1\leqslant k\leqslant n$, $B_{k}$ 的半径小于 $1/k$ 且 $\overline{B}_{k}\subseteq\mathcal{O}_{k}$. 由于 $\mathcal{O}_{n+1}$ 在 $X$ 中稠密, 集合 $B_{n}\cap\mathcal{O}_{n+1}$ 是非空的. 令 $x_{n+1}$ 属于开集 $B(x_{n},r_{n})\cap\mathcal{O}_{n}$. 选取 $r_{n+1}$, $0<r_{n+1}<1/(n+1)$, 使得若我们定义 $B_{n+1}=B(x_{n+1},r_{n+1})$, 则 $\overline{B}_{n+1}\subseteq B_{n}\cap\mathcal{O}_{n+1}$. 这归纳地定义了收缩的闭集序列 $\{\overline{B}_{n}\}_{n=1}^{\infty}$, 它具有性质: 对每个 $n$, $\overline{B}_{n}\subseteq\mathcal{O}_{n}$. 度量空间 $X$ 是完备的. 因此我们从 Cantor 交定理推出 $\bigcap_{n=1}^{\infty}\overline{B}_{n}$ 是非空的. 令 $x_{\ast}$ 属于这个交. 则 $x_{\ast}$ 属于 $\bigcap_{n=1}^{\infty}\mathcal{O}_{n}$. 另一方面, 根据\eqref{align:ch10-baire-7}, $x_{\ast}$ 也属于 $B(x_{0},r_{0})$. 这完成了(i)的证明.
\end{proof}

度量空间 $X$ 的子集 $E$ 称为无处稠密的, 若它的闭包 $\overline{E}$ 是中空的. $X$ 的子集 $E$ 是无处稠密的当且仅当对 $X$ 的每个开集 $\mathcal{O}$, $E\cap\mathcal{O}$ 在 $\mathcal{O}$ 中不稠密(见习题 16). Baire 范畴定理有以下等价的阐述: 在一个完备度量空间中, 无处稠密集的可数族的并是中空的.

\begin{corollary}\label{corollary:ch10-baire推论4}
令 $X$ 为完备度量空间, 而 $\{F_{n}\}_{n=1}^{\infty}$ 是 $X$ 的闭子集的可数族. 若 $\bigcup_{n=1}^{\infty}F_{n}$ 有非空内部, 则至少一个 $F_{n}$ 有非空内部. 特别地, 若 $X=\bigcup_{n=1}^{\infty}F_{n}$, 则至少一个 $F_{n}$ 有非空内部.
\end{corollary}

\begin{corollary}\label{corollary:ch10-baire推论5}
令 $X$ 为完备度量空间, 而 $\{F_{n}\}_{n=1}^{\infty}$ 是 $X$ 的闭子集的可数族. 则 $\bigcup_{n=1}^{\infty}\operatorname{bd}F_{n}$ 是中空的.
\end{corollary}

\begin{proof}
我们把证明 $X$ 的任何闭子集 $E$ 的边界 $\operatorname{bd}E$ 是中空的留作练习. $X$ 的任何子集的边界是闭的. 因此, 对每个自然数 $n$, $\operatorname{bd}F_{n}$ 是闭与中空的. 根据 Baire 范畴定理, $\bigcup_{n=1}^{\infty}\operatorname{bd}F_{n}$ 是中空的.
\end{proof}

\begin{theorem}\label{theorem:ch10-baire定理6}
令 $\mathcal{F}$ 为完备度量空间 $X$ 上的连续实值函数族. 它在以下意义下逐点有界: 对每个 $x\in X$, 存在常数 $M_{x}$, 使得
\begin{align*}
\text{对所有 }f\in\mathcal{F},\quad|f(x)|\leqslant M_{x}.
\end{align*}
则存在 $X$ 的非空开子集 $\mathcal{O}$ 使得 $\mathcal{F}$ 在其上是一致有界的: 存在常数 $M$, 使得在 $\mathcal{O}$ 上
\begin{align}
\text{对所有 }f\in\mathcal{F},\quad|f|\leqslant M.\label{align:ch10-baire-8}
\end{align}
\end{theorem}

\begin{proof}
对每个指标 $n$, 定义 $E_{n}=\{x\in X\mid|f(x)|\leqslant n,\text{ 对所有 }f\in\mathcal{F}\}$. 则 $E_{n}$ 是闭的, 由于 $\mathcal{F}$ 中的函数是连续的而闭集族的交是闭的. 由于 $\mathcal{F}$ 逐点有界, 对每个 $x\in X$, 存在指标 $n$ 使得对所有 $f\in\mathcal{F}$, $|f(x)|\leqslant n$, 即 $x$ 属于 $E_{n}$. 因此 $X=\bigcup_{n=1}^{\infty}E_{n}$. 由于 $X$ 是完备度量空间, 我们从推论 4 得出存在自然数 $n$ 使得 $E_{n}$ 包含开球 $B(x,r)$. 因此\eqref{align:ch10-baire-8}对 $\mathcal{O}=B(x,r)$ 和 $M=n$ 成立.
\end{proof}

我们已经看到若连续实值函数序列一致收敛, 则极限函数是连续的, 而对于逐点收敛这是错的. 然而, 对于在完备度量空间上逐点收敛于连续实值函数序列的实值函数, 极限函数在其定义域的一个稠密子集的每一点是连续的.

\begin{theorem}\label{theorem:ch10-baire定理7}
令 $X$ 为完备度量空间, 而 $\{f_{n}\}$ 是 $X$ 上的连续实值函数序列. 它在 $X$ 上逐点收敛于实值函数 $f$. 则存在 $X$ 的稠密子集 $D$, 使得 $\{f_{n}\}$ 在 $D$ 上等度连续且 $f$ 在 $D$ 的每个点连续.
\end{theorem}

\begin{proof}
令 $m$ 和 $n$ 为自然数, 定义
\begin{align*}
E(m,n)=\{x\in X\mid|f_{j}(x)-f_{k}(x)|\leqslant 1/m,\text{ 对所有 }j,k\geqslant n\}.
\end{align*}
由于每个函数 $x\mapsto|f_{j}(x)-f_{k}(x)|$ 是连续的, 集合 $E(m,n)$ 作为闭集族的交是闭的. 根据推论 5,
\begin{align*}
D=X\sim\left[\bigcup_{n,m\in\mathbf{N}}\operatorname{bd}E_{m,n}\right]
\end{align*}
在 $X$ 中稠密. 观察到若 $n$ 和 $m$ 是自然数, 且 $D$ 中的点 $x$ 属于 $E(m,n)$, 则 $x$ 属于 $E(m,n)$ 的内部. 我们宣称 $\{f_{n}\}$ 在 $D$ 的每个点等度连续. 事实上, 令 $x_{0}$ 属于 $D$. 令 $\varepsilon>0$. 选取自然数 $m$ 使得 $1/m<\varepsilon/4$. 由于 $\{f_{n}(x_{0})\}$ 收敛到实数, $\{f_{n}(x_{0})\}$ 是 Cauchy 的. 选取自然数 $N$ 使得对所有 $j,k\geqslant N$,
\begin{align}
|f_{j}(x_{0})-f_{k}(x_{0})|\leqslant 1/m.\label{align:ch10-baire-9}
\end{align}
因此 $x_{0}$ 属于 $E_{m,N}$. 如同我们上面观察到的, $x_{0}$ 属于 $E(m,N)$ 的内部. 选取 $r>0$ 使得 $B(x_{0},r)\subseteq E(m,N)$, 即
\begin{align}
\text{对所有 }j,k\geqslant N\text{ 和所有 }x\in B(x_{0},r),\quad|f_{j}(x)-f_{k}(x)|\leqslant 1/m.\label{align:ch10-baire-10}
\end{align}
函数 $f_{N}$ 在 $x_{0}$ 连续, 因此存在 $\delta$, $0<\delta<r$, 使得
\begin{align}
\text{对所有 }x\in B(x_{0},\delta),\quad|f_{N}(x)-f_{N}(x_{0})|\leqslant 1/m.\label{align:ch10-baire-11}
\end{align}
观察到对每个点 $x\in X$ 和自然数 $j$,
\begin{align*}
f_{j}(x)-f_{j}(x_{0})=[f_{j}(x)-f_{N}(x)]+[f_{N}(x)-f_{N}(x_{0})]+[f_{N}(x_{0})-f_{j}(x_{0})].
\end{align*}
我们从\eqref{align:ch10-baire-9}、\eqref{align:ch10-baire-10}、\eqref{align:ch10-baire-11}以及三角不等式推出
\begin{align}
\text{对所有 }j\geqslant N\text{ 和所有 }x\in B(x_{0},\delta),\quad|f_{j}(x)-f_{j}(x_{0})|\leqslant 3/m\leqslant[3/4]\varepsilon.\label{align:ch10-baire-12}
\end{align}
连续函数的有限族 $\{f_{j}\}_{j=1}^{N-1}$ 显然在 $x_{0}$ 等度连续. 因此我们从\eqref{align:ch10-baire-12}推出 $\{f_{n}\}$ 在 $x_{0}$ 等度连续. 这蕴涵 $f$ 在 $x_{0}$ 的连续性. 的确, 在\eqref{align:ch10-baire-12}取 $j\to\infty$ 的极限得到
\begin{align}
\text{对所有 }x\in B(x_{0},\delta),\quad|f(x)-f(x_{0})|<\varepsilon.\label{align:ch10-baire-13}
\end{align}
\end{proof}

\begin{remark}\label{remark:ch10-第一范畴}
有一个与本节思想相关的标准术语. 度量空间 $X$ 的子集 $E$ 称为第一范畴的(或疏朗的), 若 $E$ 是 $X$ 的可数个无处稠密的子集的并. 不是第一范畴的集合称为第二范畴的. Baire 范畴定理也可重述为: 完备度量空间的非空开子集是第二范畴的.
\end{remark}

\begin{remark}\label{remark:ch10-baire应用}
Baire 范畴定理的推论非常多样. 在第 13 章, 我们用定理 6 证明开映射定理和一致有界原理, 它们是研究线性泛函和算子的两个基石. 在第 18 章, 我们用定理 7 证明关于测度收敛的 Vitali-Hahn-Saks 定理, 它是描述在 $L^{1}(E)$ 中弱收敛的实质要素. 在习题 20 和 21, 从 Baire 范畴定理导出连续可微函数的两个有趣性质.
\end{remark}

\begin{exercise}\label{exercise:ch10-15}
令 $E$ 为完备度量空间 $X$ 的子集. 证明 $\operatorname{bd}E$ 是闭的. 证明: 若 $E$ 是闭的, 则 $\operatorname{bd}E$ 的内部是空的.
\end{exercise}

\begin{exercise}\label{exercise:ch10-16}
在一个度量空间 $X$ 中, 证明子集 $E$ 无处稠密当且仅当对 $X$ 的每个开子集 $\mathcal{O}$, $E\cap\mathcal{O}$ 不在 $\mathcal{O}$ 中稠密.
\end{exercise}

\begin{exercise}\label{exercise:ch10-17}
在完备度量空间 $X$ 中, 无处稠密集的可数族的并也是无处稠密的吗?
\end{exercise}

\begin{exercise}\label{exercise:ch10-18}
令 $\mathcal{O}$ 为开子集而 $F$ 为度量空间 $X$ 的闭子集. 证明 $\overline{\mathcal{O}}\sim\mathcal{O}$ 和 $F\sim\operatorname{int}F$ 是闭与中空的.
\end{exercise}

\begin{exercise}\label{exercise:ch10-19}
在完备度量空间中, 第一范畴集的可数族的并也是第一范畴的吗?
\end{exercise}

\begin{exercise}\label{exercise:ch10-20}
令 $F_{n}$ 为 $C[0,1]$ 的子集, $C[0,1]$ 由函数组成, 其中存在 $[0,1]$ 中的点 $x_{0}$, 使得对所有 $x\in[0,1]$,
\begin{align*}
|f(x)-f(x_{0})|\leqslant n|x-x_{0}|.
\end{align*}
证明 $F_{n}$ 是闭的. 通过观察到对 $f\in C[0,1]$ 与 $r>0$, 存在分段线性函数 $g\in C[0,1]$ 使得 $\rho_{\max}(f,g)<r$, 且 $g$ 在 $[0,1]$ 上的左右导数比 $n+1$ 大, 证明 $F_{n}$ 是中空的. 得出
\begin{align*}
C[0,1]\neq\bigcup_{n=1}^{\infty}F_{n}
\end{align*}
并证明每个 $h\in C[0,1]\sim\bigcup_{n=1}^{\infty}F_{n}$ 在 $(0,1)$ 中的任何点不可微.
\end{exercise}

\begin{exercise}\label{exercise:ch10-21}
令 $f$ 为度量空间 $X$ 上的实值函数. 证明使得 $f$ 在其上是连续的点集是开集的可数族的交. 得出 $\mathbf{R}$ 上不存在仅在有理点连续的实值函数.
\end{exercise}

\begin{exercise}\label{exercise:ch10-22}
对每个自然数 $n$, 证明 $[0,1]$ 中存在无处稠密的闭集, 它的 Lebesgue 测度为 $1-1/n$. 用此构造出 $[0,1]$ 中测度为 $1$ 的第一范畴集.
\end{exercise}

\begin{exercise}\label{exercise:ch10-23}
度量空间 $X$ 中的点 $x$ 称为孤立的, 若单点集 $\{x\}$ 在 $X$ 中是开的.
\begin{enumerate}[(i)]
\item 证明没有孤立点的完备度量空间有不可数个点.
\item 用(i)部分证明 $[0,1]$ 是不可数的. 将此与 $[0,1]$ 是不可数的(因为它有正的 Lebesgue 测度)的证明做比较.
\item 证明: 若 $X$ 是没有孤立点的完备度量空间, 而 $\{F_{n}\}_{n=1}^{\infty}$ 是闭中空集的可数族, 则 $X\sim\{F_{n}\}_{n=1}^{\infty}$ 是稠密与不可数的.
\end{enumerate}
\end{exercise}

\begin{exercise}\label{exercise:ch10-24}
令 $E$ 为完备度量空间 $X$ 的子集. 证明以下断言.
\begin{enumerate}[(i)]
\item 若 $X\sim E$ 是稠密的而 $F$ 是包含于 $E$ 的闭子集, 则 $F$ 无处稠密.
\item 若 $E$ 和 $X\sim E$ 都是稠密的, 则它们中至多有一个是闭集的可数族的并.
\item $[0,1]$ 中的有理数集不是开集的可数族的交.
\end{enumerate}
\end{exercise}

\begin{exercise}\label{exercise:ch10-25}
证明在定理 6 的假设下存在稠密开集 $\mathcal{O}\subseteq X$, 对每个 $x\in\mathcal{O}$ 有邻域 $U$ 使得在其上 $\mathcal{F}$ 是一致有界的.
\end{exercise}

\begin{exercise}\label{exercise:ch10-26}
根据 H\"older 不等式, 我们有 $L^{2}[a,b]\subseteq L^{1}[a,b]$. 证明 $L^{2}[a,b]$ 作为完备度量空间 $L^{1}[a,b]$ 的子集, 是第一范畴的.
\end{exercise}

\begin{exercise}\label{exercise:ch10-27}
令 $f$ 为 $\mathbf{R}$ 上的连续实值函数, 它具有性质: 对每个实数 $x$, $\lim_{n\to\infty}f(nx)=0$. 证明 $\lim_{x\to\infty}f(x)=0$.
\end{exercise}

\begin{exercise}\label{exercise:ch10-28}
令 $f$ 为 $\mathbf{R}$ 上的连续实值函数. 假定对每个实数 $x$, 存在指标 $n=n(x)$ 使得 $f^{(n)}(x)=0$. 证明 $f$ 是多项式. (提示: 用两次 Baire 范畴定理.)
\end{exercise}

\end{document}